\documentclass{article}
\usepackage[utf8]{inputenc}
\usepackage[T1]{fontenc}
\usepackage{amsmath, amssymb}
\usepackage{geometry}
\usepackage{hyperref}
\usepackage{booktabs}
\usepackage{caption}

\title{Technical Report: Are the Electron Mass and Classical Electron Radius\\[0.2em] Inputs to the Enhanced EWT Model?}
\author{Łukasz Smoliński\\[0.2em]\texttt{l\_smolinski@o2.pl}}
\date{\today}

\begin{document}
	
	\maketitle
	
	\begin{abstract}
		This report examines the role of the electron mass $m_e$ and the classical electron radius $r_e$ within the Enhanced Energy Wave Theory (EWT). It addresses the fundamental question of whether these quantities serve as empirical inputs to the model or emerge as derived consequences of a deeper geometric structure. The analysis distinguishes between an initial calibration phase, where empirical values are used to fix the absolute scale of the vacuum lattice, and the final zero‑parameter formulation, in which all major physical constants are expressed purely in terms of $\pi$, Euler's number $e$, the integers $8$ (BCC coordination) and $10$ (electron wave‑center count). It is shown that after calibration, $r_e$ and $m_e$ are no longer inputs: $r_e$ follows from the statutory neutrino radius $r_\nu$ via $r_e = 100\,r_\nu$, and $m_e$ does not appear in the derivation of the fine‑structure constant $\alpha$, the gravitational constant $G$, the lepton anomalous magnetic moments $a_l$, the atomic scales ($R_\infty$, $a_0$, $\lambda_C$), or the mixing angles $\theta_W$, $\theta_C$. The model therefore achieves genuine zero‑parameter predictive power, with all emergent quantities matching experiment to high precision.
	\end{abstract}
	
	\section{Introduction}
	A recurring question when assessing a new theoretical framework is whether its fundamental constants are truly derived from first principles or merely fitted to experimental data. In the Enhanced Energy Wave Theory (EWT) presented in the accompanying paper, quantities such as the electron mass $m_e$ and the classical electron radius $r_e$ appear early in the calculations. This report clarifies their status: they are not free parameters of the final theory but serve as temporary anchors during the calibration of the geometric scale. Once the lattice geometry is fixed, all essential constants become emergent predictions.
	
	\section{Categories of Quantities}
	To make the discussion precise, three classes of quantities are distinguished:
	
	\begin{itemize}
		\item \textbf{Fundamental geometric inputs:} pure numbers that require no experimental adjustment – $\pi$, Euler's number $e$, the integer $8$ (coordination number of the BCC lattice), and the integer $10$ (wave‑center count $K_{WC}$ for the electron).
		\item \textbf{Calibration anchors:} empirical constants used during the construction of the model to determine the absolute scale of the lattice. In this phase one may employ $m_e$, $r_e$, or the Planck charge $q_P$, but after the scale is fixed these anchors are replaced by geometric relations.
		\item \textbf{Emergent derived quantities:} constants that can be calculated from the geometric inputs alone, requiring no further empirical information. This class includes $\alpha$, $G$, $a_e$, $a_\mu$, $a_\tau$, $R_\infty$, $a_0$, $\lambda_C$, $\theta_W$, $\theta_C$, and others.
	\end{itemize}
	
	\section{The Calibration Phase: Fixing the Lattice Scale}
	Every physical theory must connect its abstract concepts to measurable quantities. In EWT the fundamental length scale is the statutory neutrino radius $r_\nu$. It is obtained from the Planck charge $q_P$ and Euler's number $e$:
	\begin{equation}
		r_\nu = \frac{2 q_P e^2}{g_v},
	\end{equation}
	where $g_v$ is a geometric correction factor. The value of $q_P$ is taken from CODATA, but within EWT it is reinterpreted as a wave amplitude measured in meters – a step already present in Yee's earlier work \cite{yee2020constants}. Thus $q_P$ acts as a calibration anchor, not as a free parameter of the final theory.
	
	\subsection{The Role of the Geometric Factor $g_v$}
	The geometric correction factor $g_v$ is applied to maintain numerical consistency with the core EWT model. While the fundamental Lorentz-like shortening of matter in motion is an integral part of the EWT framework, the specific factor $g_v$ used to define $r_\nu$ (which results in radius expansion relative to the uncorrected wavelength) is not interpreted as a kinematic Doppler shift. Instead, it is better interpreted as a consequence of the absence, or near-absence, of the magnetic deformation term $\epsilon_M$ in the neutral neutrino. This term is critical for describing the electron's spin and anomalous magnetic moment, suggesting the expansion is driven by the fundamental difference in magnetic/spin geometry between the two leptons.
	
	Specifically, it is proposed that the magnetic field of a charged lepton acts as a geometric torque that actively compresses the soliton radius $r$, whereas the absence of this strain in the neutral neutrino allows it to maintain its expanded statutory radius $r_\nu$.
	
	The numerical consistency script (see Appendix) verifies that this statutory value $r_\nu$ adheres to the power-law relationships that govern the ratios of lepton soliton radii. The numerical test confirms that $r_\nu$ yields an implied geometric scaling factor of $\approx 10^{10}$, validating its derived magnitude with exceptional precision (see Part II of the output). The model's robustness against large relative changes in radius is further demonstrated in the accompanying sensitivity analysis.
	
	\section{Validation of the 1:100 Resonance}
	Using the empirical $r_e$, the model then verifies the predicted 1:100 resonance:
	\begin{equation}
		\frac{r_e}{r_\nu} = 100.000021,
	\end{equation}
	confirming that $r_e = 100\,r_\nu$ holds with extraordinary precision. This relation is not an assumption; it is a successful postdiction that validates the choice $K_{WC}=10$ for the electron.
	
	\section{The Fine‑Structure Constant: A Purely Geometric Prediction}
	One of the most important results of EWT is the derivation of the fine‑structure constant from pure geometry:
	\begin{equation}
		\alpha^{-1} = (4\pi^3 + \pi^2 + \pi) - \epsilon_M,\qquad \epsilon_M = \frac{1}{8\pi^7}.
	\end{equation}
	This expression contains only $\pi$ and the integer $8$ (the BCC coordination number). The numerical result $\alpha^{-1}=137.036262$ matches the CODATA value $137.035999$ to within $2\times10^{-4}\%$. \textbf{Once the geometric framework is established, $\alpha$ is a prediction, not an input.}
	
	\section{Emergent Quantities: What the Model Actually Predicts}
	After calibration, the model uses only $r_\nu$ and $\alpha$ to derive all other constants. The electron radius itself becomes an emergent quantity:
	\begin{equation}
		r_e = 100\,r_\nu.
	\end{equation}
	From $r_e$ and $\alpha$ the three fundamental atomic scales follow directly:
	
	\begin{align}
		R_\infty &= \frac{\alpha^3}{4\pi r_e} = \frac{\alpha^3}{400\pi r_\nu},\\[2mm]
		a_0     &= \frac{r_e}{\alpha^2} = \frac{100\,r_\nu}{\alpha^2},\\[2mm]
		\lambda_C &= \frac{2\pi r_e}{\alpha} = \frac{200\pi r_\nu}{\alpha}.
	\end{align}
	
	Numerical evaluation using the geometric values yields:
	
	\[
	\begin{array}{lccc}
		\text{Constant} & \text{EWT Prediction} & \text{CODATA 2022} & \text{Error (ppm)}\\\hline
		R_\infty\;(\text{m}^{-1}) & 10973670.62263460 & 10973731.56815700 & 5.55\\
		a_0\;(\text{m}) & 5.291791330634349\times10^{-11} & 5.291772109030000\times10^{-11} & 3.63\\
		\lambda_C\;(\text{m}) & 2.426314389900505\times10^{-12} & 2.426310238670000\times10^{-12} & 1.71
	\end{array}
	\]
	
	All three emerge from the same two geometric inputs – $r_\nu$ and $\alpha$ – demonstrating internal consistency.
	
	\section{The Electron Mass $m_e$ and Lepton Anomalous Magnetic Moments}
	The electron mass is not required as an input for any of the derivations listed above. It appears only in the definition of the base gravitational constant $G_{\text{Base}} = c^2 r_e/m_e$, which is used to construct the full gravitational operator $\mathcal{U}$. In the zero‑parameter formulation, $m_e$ itself can be obtained from $r_e$ and $\alpha$ via the relation $r_e = \alpha\hbar/(m_e c)$, but this is a consistency check rather than a necessary input.
	
	\textbf{Crucially, the calculation of the anomalous magnetic moments of the leptons ($a_e$, $a_\mu$, $a_\tau$) does not require the electron mass.} The electron AMM $a_e$ is derived directly from pure geometry in Part XI of the script:
	
	\begin{equation}
		a_e = \frac{8\pi^4 - 1}{64\pi^8 + 16\pi^7 + 16\pi^6 - 2\pi^{-2}},
	\end{equation}
	
	an expression containing no empirical inputs whatsoever. The muon and tau AMMs follow recursively from the same geometric modulator $\epsilon_M$, with the Fibonacci invariants $L_\mu = 5$ and $L_\tau = 34$ acting as structural latches. The numerical results are:
	
	\[
	\begin{array}{lccc}
		\text{Generation} & \text{EWT Prediction} & \text{Experimental Target} & \text{Relative Error}\\\hline
		\text{Electron } a_e & 1.159917127722\times10^{-3} & 1.159652181600\times10^{-3} & 0.0228\%\\
		\text{Muon } a_\mu & 248.572\times10^{-6} & 248.800\times10^{-6} & 0.091\%\\
		\text{Tau } a_\tau & 1176.845\times10^{-6} & 1177.210\times10^{-6} & 0.031\%
	\end{array}
	\]
	
	These predictions are completely independent of the electron mass, confirming that lepton anomalies are purely geometric properties of the BCC lattice.
	
	\section{Other Emergent Quantities}
	The same geometric logic produces the gravitational constant $G$ and the mixing angles $\theta_W$, $\theta_C$. In every case the final expressions contain only $\pi$, $e$, and the integers $8$, $10$, $5$, $34$ – no free parameters remain.
	
	\section{Conclusion}
	In the final, zero‑parameter formulation of Enhanced EWT, \textbf{neither the electron mass $m_e$ nor the classical electron radius $r_e$ serve as inputs}. They are emergent quantities:
	\begin{itemize}
		\item $r_e$ is fixed by the statutory neutrino radius $r_\nu$ through $r_e = 100\,r_\nu$, with $g_v$ accounting for the magnetic deformation difference between charged and neutral leptons;
		\item $m_e$ does not appear in the derivation of any major constant; it is irrelevant for predicting $\alpha$, $G$, the lepton anomalies, the atomic scales, or the mixing angles;
		\item the anomalous magnetic moments $a_e$, $a_\mu$, $a_\tau$ are calculated without any reference to $m_e$, demonstrating that lepton anomalies are pure geometric properties of the vacuum lattice.
	\end{itemize}
	The model's predictive power is demonstrated by the fact that more than a dozen physical constants are derived from a handful of purely geometric numbers, with all predictions lying within $0.1\%$ of experiment and many reaching sub‑ppm accuracy. This constitutes a genuine paradigm shift from empirical parameter‑fitting to structural determinism.
	
	\begin{thebibliography}{9}
		\bibitem{yee2020constants} Yee, J. (2020). Fundamental Physical Constants: Explained and Derived by Energy Wave Equations. \textit{Preprint}.
	\end{thebibliography}
	
\end{document}